\subsection{A rational symmetry reduction of the cubic action class}

\begin{proposition}[E.23 --- Proved: two projected arithmetic parameters and an eight-dimensional preimage]
\label{E-cubic-symmetry}
Retain the actual cubic notation of E.18--E.19. Put
$H=T_3E$, where $E=X_3$, and write
$V_{\mathbb Q_3}=H_{\mathbb Q_3}\oplus W_{\mathbb Q_3}$
for its rational compact and boundary summands. Let
$w\in\bigwedge^2V$ be the primitive compact determinant
vector of E.19, and define the actual submodule and quotient
\[
 B_w=[V,w]\subset A=\operatorname{Lie}_3(V),
 \qquad \mathcal Q=A/B_w.
 \tag{E23.1}
\]
Then $B_w$ is primitive of rank ten, its bracket map identifies
it with $V(1)$ as a $G_{\mathbb Q}$-module, and
$\mathcal Q$ is free of rank $320$.

Let $K=\mathbb Q(\zeta_3)$ and let
$\mathcal A=E[3]\rtimes\mu_6$ be the order-$54$ group of
$K$-automorphisms of $U=E\setminus E[3]$. Its action on
$\mathcal Q$ is independent of base paths. The subspace
\[
 D=(\mathcal Q_{\mathbb Q_3})^{\mathcal A}
\]
is $G_{\mathbb Q}$-stable and has the following rational
$G_{\mathbb Q}$-module description:
\[
 D\simeq H_{\mathbb Q_3}(2)\oplus
        \mathbb Q_3(3)\oplus
        \mathbb Q_3(3)\varepsilon_K,
 \tag{E23.2}
\]
where $\varepsilon_K$ is the quadratic character of $K/\mathbb Q$.
In particular $\dim D=4$ and $\operatorname{tr}(c|D)=0$.

For $\mathcal G=G_{\mathbb Q,\{3,\infty\}}$, the image
$\overline\eta$ of the actual class of E.19 satisfies
\[
 \overline\eta\in
 H^1_{\rm cts}(\mathcal G,D(-1))
 \lhook\joinrel\longrightarrow
 H^1_{\rm cts}(\mathcal G,\mathcal Q_{\mathbb Q_3}(-1)),
 \qquad
 \dim_{\mathbb Q_3}H^1(\mathcal G,D(-1))=2.
 \tag{E23.3}
\]
The two summands with potentially nonzero classes in (E23.2)
after twisting by $(-1)$ are $H_{\mathbb Q_3}(1)$ and
$\mathbb Q_3(2)\varepsilon_K$; each has one-dimensional
unramified global $H^1$. The untwisted quadratic-character
summand $\mathbb Q_3(2)$ has zero $H^1$.

More explicitly, the original rational class $\eta$ lies in
an at-most-eight-dimensional subspace of
$H^1(\mathcal G,A_{\mathbb Q_3}(-1))$: it is the inverse
image of the two-dimensional subspace in (E23.3) under the
map in the exact sequence
\[
 H^1(\mathcal G,V_{\mathbb Q_3})\longrightarrow
 H^1(\mathcal G,A_{\mathbb Q_3}(-1))\longrightarrow
 H^1(\mathcal G,\mathcal Q_{\mathbb Q_3}(-1)),
 \qquad \dim H^1(\mathcal G,V_{\mathbb Q_3})=6.
 \tag{E23.4}
\]
These are constraints on the actual arithmetic class and the
actual quotient maps. They do not evaluate its two projected
coordinates, its component in the displayed kernel, or its
integral torsion components. No identification of those
coordinates with specified polylogarithmic values is asserted.
\end{proposition}
\begin{proof}
Write $w=u\wedge v$, with
$u=3f_1+U_1$ and $v=3f_2+U_2$ as in E.19. A coefficient
of $w$ is a unit. Equivalently a $2\times2$ minor of the
matrix $(u,v)$ is a unit, so $u,v$ extend to a
$\mathbb Z_3$-basis of $V$. In that basis the ten elements
$[e_i,[u,v]]$ are independent in the degree-three free Lie
module. For $e_i=u,v$ they are two individual Hall basis
elements up to sign. For each remaining $e_i$, the component
of multidegree $u v e_i$ has rank two, and
$[e_i,[u,v]]$ is a vector with a unit coordinate there.
Different indices have disjoint multidegrees. These unit
coordinates provide an integral left inverse on their span;
thus the span is primitive. Galois equivariance of the bracket
and $g(w)=\chi(g)w$ identify the span with $V(1)$.
This proves (E23.1), including its integral assertion.

Let $L_w$ be the inverse image of $\mathbb Z_3w$ under
$L=\gamma_2N/\gamma_4N\to\bigwedge^2V$.
An inner automorphism by an element of $N$ with image
$a\in V$ acts trivially on $A$, and sends a lift of
$b w$ to that lift plus $b[a,w]$. Consequently it acts
trivially on the full extension
\[
 0\longrightarrow\mathcal Q\longrightarrow
 L_w/B_w\longrightarrow\mathbb Z_3(1)w
             \longrightarrow0.
 \tag{E23.5}
\]
This is the exact quotient that removes the base-path
ambiguity. In particular an automorphism of the curve
preserving $w$ acts genuinely on (E23.5), although its
action on $N$ was originally specified only up to inner
conjugation. Commutators of lifts of commuting outer
automorphisms are inner, hence disappear on (E23.5) as well.

We check preservation of the compact line used here.
The generalized-Jacobian splitting $j_0$, with proper
projection $p j_0=[3]$, is unique because a homomorphism
from an abelian variety to a torus is zero. It therefore
commutes with the automorphisms of the pair $(E,E[3])$.
Translations act identically on its compact part; the
units in $\mu_6$ have determinant one on $T_3E$.
Thus $\mathcal A$ fixes $w$ literally. Its automorphisms
are defined over $K$ and are permuted by
$\operatorname{Gal}(K/\mathbb Q)$. The preceding inner
calculation proves that its action on (E23.5) commutes
with $G_K$ and is normalized by the entire $G_{\mathbb Q}$.
No chosen path has been declared Galois invariant.

Over $\mathbb Q_3$, average any lift of $w$ in (E23.5)
over $\mathcal A$. This is legitimate because $54$ is
invertible in $\mathbb Q_3$. It gives an
$\mathcal A$-fixed lift $\omega_{\mathcal A}$.
Replace it by
$(\omega_{\mathcal A}-c\omega_{\mathcal A})/2$.
Complex conjugation normalizes $\mathcal A$, so the
new lift is still fixed by $\mathcal A$ and now satisfies
$c\omega_{\mathcal A}=-\omega_{\mathcal A}$.
For every arithmetic $g$, both $g\omega_{\mathcal A}$
and $\chi(g)\omega_{\mathcal A}$ are
$\mathcal A$-fixed. Their difference, divided by $\chi(g)$,
is consequently a cocycle valued in $D(-1)$.
Changing the lift adds a coboundary, as in E.19, so it
represents exactly the projected class $\overline\eta$.
The averaging idempotent is $G_{\mathbb Q}$-equivariant
because Galois permutes the elements of $\mathcal A$.
It splits the inclusion $D\subset\mathcal Q_{\mathbb Q_3}$,
so the cohomology arrow in (E23.3) is injective.

For completeness, we give the finite character calculation
and its arithmetic interpretation. On the nine cusp points
identify $E[3]$ with $\mathbb F_3^2$. A generator of
$\mu_6$ acts by $-R$, where $R$ is the matrix in E18.9,
and complex conjugation acts by $C_E$. An element
$(t,\alpha)$ acts on the cusp set by $x\mapsto t+\alpha x$.
Its trace on $W_{\mathbb Q_3}$ is the number of fixed
cusp points minus one. Its trace on $H_{\mathbb Q_3}$
is $\alpha+\alpha^{-1}$, interpreted in the two-dimensional
CM representation. For a complex-conjugation coset this
compact trace is zero, and the boundary trace has an
additional minus sign from the Tate twist.
For any such operator $b$, the degree-three trace is
\[
 \operatorname{tr}(b|\operatorname{Lie}_3V)
 =\frac{\operatorname{tr}(b|V)^3-
                    \operatorname{tr}(b^3|V)}{3}.
 \tag{E23.6}
\]
Subtract $\operatorname{tr}(b|V)\det(b|H)$ for $B_w$.
Averaging over the $54$ affine transformations gives
\[
 \dim V^{\mathcal A}=0,\quad
 \dim A_{\mathbb Q_3}^{\mathcal A}=4,\quad
 \dim D=4,\quad \operatorname{tr}(c|D)=0.
\]
Keeping the number of compact and boundary factors gives
\[
\begin{array}{c|rrrr}
 \text{multidegree in }\mathcal Q&H^3&H^2W&HW^2&W^3\\\hline
 \text{invariant dimension}&0&0&2&2\\
 \text{trace of }c&0&0&0&0 .
\end{array}
 \tag{E23.7}
\]
The executed program \texttt{class\_three\_symmetry.py}
records every character row using integer arithmetic.
The independent character check accompanying this proof
repeats the $54$ affine transformations and their conjugation
coset, and checks each multidegree separately. It also records
the correction to an initial guess that all four invariant
coordinates had multidegree $HW^2$; (E23.7) is the checked
replacement.

All cusps are rational over $K$, so $G_K$ acts on $W$ by
the scalar $\chi$. The two $HW^2$ invariants couple the
two distinct CM eigenlines of $H$ to their inverse
$\mu_6$-characters. Each coupling has multiplicity one:
the nonzero characters of the translation group have one
$\mu_6$-orbit of length six and one of length two; only
the length-six orbit supplies the two primitive order-six
characters of $H$. Thus their restrictions to $G_K$ are
exactly the two CM characters times $\chi^2$.
Complex conjugation exchanges the two lines, identifying
this representation over $\mathbb Q_3$ with
$H_{\mathbb Q_3}(2)$. This can also be checked after the
quadratic coefficient extension splitting the CM algebra;
changing its exchange convention gives the quadratic twist,
which is isomorphic by the CM endomorphism on which
complex conjugation acts as minus one.
The $W^3$ invariants restrict to two copies of $\chi^3$
on $G_K$; their zero conjugation trace gives one copy of
$\chi^3$ and one of $\chi^3\varepsilon_K$. This proves
(E23.2), with the actual restriction and descent maps.

Finally all these modules are unramified outside three,
and after twisting by $(-1)$ their good-Frobenius eigenvalues
have positive weights. Their invariant vectors are zero.
The free-pro-three arithmetic source and the order-two
averaging of E.16 give zero $H^2$. The global Euler
characteristic therefore makes $\dim H^1$ the dimension
of the minus eigenspace of complex conjugation.
For $H(1)$ this is one; for $\mathbb Q_3(2)$ it is zero;
for $\mathbb Q_3(2)\varepsilon_K$ it is one. This proves
(E23.3). On $V$, the compact trace is zero and the boundary
trace is $-(3-1)=-2$, since three cusps are real.
Its dimension is ten, so the same calculation gives
$\dim H^1(\mathcal G,V_{\mathbb Q_3})=6$.
The long exact sequence of (E23.1), twisted by $(-1)$,
gives (E23.4). Its kernel is the image of that
six-dimensional space, so the inverse image of the
specified two-dimensional subspace has dimension at most
eight. This proves the announced constraint without
assigning a value to the arithmetic class. The manuscript
anchor remains \S51.16, folios 733v--736v, scans 172--178,
and task LM7-E.23(i).
\end{proof}
